\subsection{Flat modules and torsion products}

THEOREM 2. --- Let E be a right A-module. The following conditions are equivalent:

\begin{enumerate}
\def\labelenumi{(\roman{enumi})}
\item
  E is flat;
\item
  for every left A-module F, and every integer n \textgreater{} 0, one has
\end{enumerate}

\[
\operatorname{Tor} _ {n} ^ {\mathbf {A}} (\mathbf {E}, \mathbf {F}) = 0;
\]

\begin{enumerate}
\def\labelenumi{(\roman{enumi})}
\setcounter{enumi}{2}
\tightlist
\item
  for every monogenic left A-module of finite presentation F, one has
\end{enumerate}

\[
\operatorname{Tor} _ {1} ^ {\mathbf {A}} (\mathbf {E}, \mathbf {F}) = 0;
\]

\begin{enumerate}
\def\labelenumi{(\roman{enumi})}
\setcounter{enumi}{3}
\item
  for every finitely generated left ideal a of A, the canonical map E ⊗\_A a → E is injective;
\item
  for every exact sequence of right A-modules, of the form
\end{enumerate}

\[
0 \rightarrow \mathbf {G} \xrightarrow {v} \mathbf {H} \xrightarrow {w} \mathbf {E} \rightarrow 0,
\]

and every left A-module F, the sequence

\[
0 \longrightarrow \mathrm{G} \otimes_ {\mathrm{A}} \mathrm{F} \xrightarrow {v \otimes 1} \mathrm{H} \otimes_ {\mathrm{A}} \mathrm{F} \xrightarrow {w \otimes 1} \mathrm{E} \otimes_ {\mathrm{A}} \mathrm{F} \longrightarrow 0
\]

is exact.

\begin{enumerate}
\def\labelenumi{(\roman{enumi})}
\item
  \(\Rightarrow\) (ii): this is the cor. to prop. 5 of X, p. 69.
\item
  \(\Rightarrow\) (iii): this is trivial.
\item
  \(\Leftrightarrow\) (iv): every monogenic left A-module of finite presentation is isomorphic to a quotient A/a, where a is a finitely generated left ideal, so that (iii) is equivalent to (iv) by X, p. 72, example.
\item
  \(\Rightarrow\) (i): by X, p. 8, prop. 3, X, p. 72, th. 1, E is flat as soon as \(\operatorname{Tor}_{1}^{\mathbf{A}}(\mathbf{E},\mathbf{F}) = 0\) for every left A-module F. If (iii) is satisfied, this is so as soon as F is monogenic and of finite presentation. By X, p. 11, prop. 7, every A-module (resp. every monogenic A-module) is a filtered inductive limit of modules of finite presentation (resp. of monogenic modules of finite presentation); one therefore sees, by X, p. 70, prop. 8, that it suffices to prove that if \(\operatorname{Tor}_{1}^{\mathbf{A}}(\mathbf{E},\mathbf{F}) = 0\) when F is monogenic, then this is so as soon as F is finitely generated. Let us therefore reason by induction on the cardinality of a generating family \((f_{1},\ldots,f_{n})\) of F; the exact sequence
\end{enumerate}

\[
0 \rightarrow \mathrm{A} f _ {1} \rightarrow \mathrm{F} \rightarrow \mathrm{F} / \mathrm{A} f _ {1} \rightarrow 0
\]

gives rise to an exact sequence

\[
\operatorname{Tor} _ {1} ^ {\mathbf {A}} (\mathrm{E}, \mathrm{A} f _ {1}) \rightarrow \operatorname{Tor} _ {1} ^ {\mathbf {A}} (\mathrm{E}, \mathrm{F}) \rightarrow \operatorname{Tor} _ {1} ^ {\mathbf {A}} (\mathrm{E}, \mathrm{F} / \mathrm{A} f _ {1}),
\]

so that \(\operatorname{Tor}_{1}^{\mathbf{A}}(\mathbf{E},\mathbf{F})=0\) since \(\operatorname{Tor}_{1}^{\mathbf{A}}(\mathbf{E},\mathbf{A}f_{1})=0\) and \(\operatorname{Tor}_{1}^{\mathbf{A}}(\mathbf{E},\mathbf{F}/\mathbf{A}f_{1})=0\) by the induction hypothesis.

\begin{enumerate}
\def\labelenumi{(\roman{enumi})}
\item
  \(\Rightarrow\) (v): this is cor. 2 to th. 1 (X, p. 72).
\item
  \(\Rightarrow\) (iii): the exact sequence (X, p. 50)
\end{enumerate}

\[
0 \longrightarrow Z _ {0} (\mathrm{E}) \stackrel {i _ {\mathrm{E}}} {\longrightarrow} L _ {0} (\mathrm{E}) \stackrel {p _ {\mathrm{E}}} {\longrightarrow} \mathrm{E} \longrightarrow 0
\]

gives rise, for every left A-module F, to an exact sequence

\[
0 \longrightarrow \operatorname{Tor} _ {1} ^ {\mathrm{A}} (\mathrm{E}, \mathrm{F}) \longrightarrow Z _ {0} (\mathrm{E}) \otimes_ {\mathrm{A}} \mathrm{F} \xrightarrow {i _ {\mathrm{E}} \otimes 1} \mathrm{L} _ {0} (\mathrm{E}) \otimes_ {\mathrm{A}} \mathrm{F} \xrightarrow {p _ {\mathrm{E}} \otimes 1} \mathrm{E} \otimes_ {\mathrm{A}} \mathrm{F} \longrightarrow 0.
\]

If (v) is satisfied, one has \(\operatorname{Tor}_1^{\mathbf{A}}(\mathbf{E},\mathbf{F}) = 0\) whence (iii).

COROLLARY 1. --- Let \(0 \to E' \to E \to E'' \to 0\) be an exact sequence of right A-modules. Suppose that \(E''\) is flat. Then for E to be flat, it is necessary and sufficient that \(E'\) be flat.

Let F be a left A-module. Since \(\operatorname{Tor}_i^{\mathbf{A}}(\mathbf{E}^{\prime \prime},\mathbf{F}) = 0\) for \(i = 1,2\) (th. 2, (i) \(\Rightarrow\) (ii)), one has an exact sequence

\[
0 \to \mathrm{Tor} _ {1} ^ {\mathbf {A}} (\mathbf {E} ^ {\prime}, \mathbf {F}) \to \mathrm{Tor} _ {1} ^ {\mathbf {A}} (\mathbf {E}, \mathbf {F}) \to 0
\]

whence the assertion (th. 2, (i) ⇔ (iii)).

COROLLARY 2. --- Let \(0 \to E_n \to E_{n-1} \to \cdots \to E_1 \to 0\) be an exact sequence of right A-modules. If \(E_i\) is flat for \(i = 1, \ldots, n-1\) , then \(E_n\) is flat.
% semantic-fragments: legacy-input-seam
\relax{}
